\section{Extending small-field derivative estimates to the boundary}
\label{sec:healthy-boundary-jets}
A small-field configuration may have a plaquette exactly at the
health threshold. Identities first proved where every plaquette is
strictly below that threshold must therefore be extended to this
boundary. We follow the actual radial argument. It uses the smallness
of the logarithmic field and a second-derivative estimate; it does
not assume that the healthy set is convex.

\subsection{A radial retreat lemma}
Let $X:[0,1]\to E$ be a twice continuously differentiable curve in
a finite-dimensional normed space, with
\[
 \sup_{0\le s\le1}\|X''(s)\|\le K<t,
             \qquad \|X(1)-X(0)\|\le t.
\]
Then $\|X(s)-X(0)\|<t$ for all $s<1$ sufficiently close to one.
If the endpoint inequality is strict, this is continuity. In the
remaining case put $D=X(1)-X(0)$, so $\|D\|=t$, and
$g(s)=X(s)-sD$. The fundamental theorem of calculus gives
\[
 \|X'(1)-D\|
  =\left\|\int_0^1\{X'(1)-X'(r)\}\,dr\right\|
                      \le K<t.
\]
Since $g'(1)=X'(1)-D$, the one-sided difference quotient satisfies
$\|g(s)-g(1)\|<(1-s)t$ near one. For $0\le s<1$ in that neighborhood,
\[
 X(s)-X(0)=sD+g(s)-g(1),\qquad
 \|X(s)-X(0)\|<st+(1-s)t=t.
\]
The inequality is strict even when the original endpoint lies exactly
on the boundary.

Let $W$ be the actual logarithmic coordinate of a coarse field, and
write $U_s=\operatorname{actual}(sW)$ for the field obtained by
exponentiating its bond coordinates. For one plaquette let
$X(s)=\rho(\operatorname{Hol}_{\square}U_s)$, where $\rho$ is the
fixed unitary representation. Then $X(0)=1$ and
\[
 \operatorname{chord}(\operatorname{Hol}_{\square}U_s)
                     =\|X(s)-1\|/r_G,
\]
with the fixed representation normalization $r_G>0$.
On a fixed small logarithmic ball the exponential has bounded first
and second derivatives. If the coordinate norm is at most $w\le1$,
its bond derivative scale is at most $Cw$. The product rule over the
four plaquette bonds gives
\[
                  \sup_{s\in[0,1]}\|X''(s)\|
                           \le r_G(4r_GCw)^2.
\]
Thus, if $(4r_GCw)^2<\tau$ and all endpoint plaquette chords are at
most $\tau$, the retreat lemma with $t=r_G\tau$ makes every plaquette
strictly healthy for $s<1$ close enough to one. There are finitely
many plaquettes, so one neighborhood works for all of them.

In the construction the logarithmic radius is
$h=\beta_F^{-7/16}$, while the health threshold has scale
$\tau=\beta_F^{-1/2}$ times its prescribed slowly varying factor.
The coordinate norm conversion has a fixed constant. Since
$h^2$ has exponent $-7/8<-1/2$, the displayed strict curvature
inequality holds eventually, uniformly over every field in that
logarithmic ball. This is the needed strict radial approximation.

\subsection{Transferring chart jets to ordered group derivatives}
We record the normalization used with this approximation. Let $j$
be fixed. Compactness on a small Lie-algebra chart supplies constants
$J,d>0$ for derivatives through order $j$ of the local logarithm.
Shrink the chart so that the logarithm inverts the exponential there.
Let $U=\operatorname{actual}(W)$, with $\|W\|\le h$, and suppose
\[
            f\circ\operatorname{actual}\in
                       \mathcal B_j(T,h^{-1};W).
\]
Consider an admissible group-valued variation $V(y)$ through $U$,
whose link derivative scale is $C$. The local logarithm and the
inverse linear coordinate map produce a chart variation $a(y)$
with $a(y_0)=W$ and $V=\operatorname{actual}\circ a$ near $y_0$.
The logarithm chain rule bounds its positive-order derivatives by
\[
             j!J\{(1+r_G)C\}^{r}\qquad(1\le r\le j).
\]
Assume
\[
       (1+r_G)C\le1,
             \qquad j!J(1+r_G)C\le h.
\]
Then every derivative of $a$ through order $j$, including its value,
is at most $h$. In the chain rule for $f\circ\operatorname{actual}\circ a$,
an outer derivative of order $m$ costs $Th^{-m}$ and its $m$
inner factors cost at most $h^m$. They cancel. Bounding the finite
partition count conservatively gives
\begin{equation}\label{eq:healthy-chart-ticket}
                       \|f\circ V\|_{j,y_0}\le j!3^jT.
\end{equation}
The same estimate applies to ordered left derivatives along any
list of at most $j$ directions of norm at most $\tau$. Indeed use
the corresponding ordered products of exponentials as the variation;
the compact exponential bounds give a link scale at most $C_jj\tau$.
Choose $C$ to dominate that value. In the present construction
$C$ has upper exponent $-1/2$, strictly below the radius exponent
$-7/16$, so both normalization inequalities hold eventually.
The construction also proves the required local differentiability
of $f$ in the group chart.

\subsection{Subtracting a summable activity array}
Let $u_i$ be regular activities with nonempty assigned supports and
nonnegative jet prices $k_i$. Suppose their rooted support-weighted
price is at most $b$ per torus site and their derivatives in directions
bounded by $\tau$ are at most $k_i$. Then
\[
                     \sum_i k_i\le |\Lambda|b.
\]
This follows by summing over the roots and removing weights at least
one. On the open set $\mathcal S$ of strictly healthy configurations,
every $u_i$ is differentiable to order $j$. Their series can be
differentiated term by term there. To check the needed domination
for any fixed list of directions, scale every direction by a common
positive number so that its norm is at most $\tau$. Multilinearity
of the ordered derivative then multiplies the summable majorant
$k_i$ by only one fixed power of that number. The common majorant
is uniform on $\mathcal S$. The local differentiation theorem for
normally convergent derivative series therefore applies, order by
order. In particular, for $U\in\mathcal S$ and normalized directions,
\begin{equation}\label{eq:healthy-array-sum}
 \left|\sum_i u_i(U)\right|\le|\Lambda|b,\qquad
 \left\|D_\ell\sum_i u_i(U)\right\|\le|\Lambda|b
                                  \quad(|\ell|\le j).
\end{equation}

Suppose the local action identity on $\mathcal S$ is
\[
                      f(U)=z-H(U)-\sum_i u_i(U),
\]
where $z$ is constant. For any function $q$ write
$\widetilde q(U)=q(U)-q(1)$. The identity field belongs to
$\mathcal S$, because $\tau>0$. Subtracting its value yields
\[
              \widetilde H(U)
             =\sum_i u_i(1)-\sum_i u_i(U)-\widetilde f(U).
\]
Assume now that the chart bound with amplitude $T$ and scale $h^{-1}$
holds for $\widetilde f\circ\operatorname{actual}$ at every $W' $
such that $\|W'\|+2d<h$ and $\operatorname{actual}(W')\in\mathcal S$.
This is a bound on the centered function itself. Apply
\eqref{eq:healthy-chart-ticket} to $\widetilde f$. Combining it with
\eqref{eq:healthy-array-sum} gives
\begin{equation}\label{eq:healthy-interior-remainder}
          \|D_\ell\widetilde H(U)\|
                   \le j!3^jT+2|\Lambda|b.
\end{equation}
The extra factor two covers the centered value at derivative order
zero; for positive orders the identity-field constant differentiates
to zero.

Now let $U=\operatorname{actual}(W)$ have no flagged plaquettes,
allowing equality at the threshold, and assume
$\|W\|+2d<h$ with $d>0$. For $0\le s\le1$,
$\|sW\|+2d<h$ as well. The radial retreat places $U_s$ in
$\mathcal S$ for $s<1$ sufficiently close to one, so the same
interior estimate applies there. Assume $H$ is regular on the full
healthy domain, as supplied by the local action construction. Its
ordered derivatives through order $j$ are continuous at $U$ in the
local group chart. Since $U_s\to U$, passing to the limit in
\eqref{eq:healthy-interior-remainder} gives the same bound at $U$.
Only $H$ needs this boundary regularity: no boundary smoothness is
inferred for the original masked series or for $f$ from its interior
identity.
